% ==========================================================================
%  Chapter 67 — Reinforced Sub-Model Integration
% ==========================================================================

\chapter{Reinforced Sub-Model Integration}
\label{ch:reinforced-submodel}

Chapters~\ref{ch:ko-bathe-3d} and~\ref{ch:truss-element} developed the two
key ingredients for modelling reinforced-concrete prismatic sub-models:
the Ko--Bathe three-dimensional concrete constitutive model and the
\texttt{TrussElement<Dim>} embedded rebar element.  This chapter describes
how these are assembled into a unified sub-model solver via
\texttt{MultiElementPolicy}, the reaction-force-based effective-modulus
extraction, and a critical bug fix in the incremental analysis engine
that affected all displacement-controlled nonlinear analyses with more
than one load step.


% =========================================================================
\section{Heterogeneous Element Assembly}
\label{sec:reinf-hetero}

A reinforced sub-model contains two distinct element types:

\begin{enumerate}
  \item \textbf{Hex8 continuum elements}
        (\texttt{ContinuumElement<ThreeDimensionalMaterial, 3, SmallStrain>})
        with the Ko--Bathe concrete law, discretising the concrete volume.
  \item \textbf{Line2 truss elements}
        (\texttt{TrussElement<3>}) with a Menegotto--Pinto steel law,
        discretising longitudinal rebars.
\end{enumerate}

Both element families share the same \texttt{Domain<3>} and the same DMPlex
sieve.  Truss elements connect pairs of hex nodes along the $z$-axis using
\texttt{Line2} topology; because their nodes coincide with hex vertices, the
perfect-bond assumption is enforced automatically through the shared
degrees of freedom.

The heterogeneous collection is stored in a
\texttt{std::vector<FEM\_Element>} and managed by \texttt{MultiElementPolicy},
which erases element types behind the \texttt{FiniteElement} interface.
The model is constructed with Constructor~2:

\begin{lstlisting}[language=C++, basicstyle=\ttfamily\small]
Model<Policy, SmallStrain, 3, MultiElementPolicy>
    M{sub.domain, std::move(all_elements)};
\end{lstlisting}


% =========================================================================
\section{Domain Builder}
\label{sec:reinf-domain-builder}

The helper \texttt{make\_reinforced\_prismatic\_domain()} in
\texttt{PrismaticDomainBuilder.hh} extends the hex-only prismatic builder
with embedded rebar:

\begin{enumerate}
  \item Build the base prismatic hex mesh (nodes and \texttt{Hex8} cells).
  \item For each bar in the \texttt{RebarSpec}, compute the grid indices
        $(i_x, i_y)$ closest to the bar's cross-section position.
  \item Create one \texttt{Line2} element per rebar segment between
        consecutive $z$-layers, reusing existing node IDs.
  \item Return a \texttt{ReinforcedDomainResult} containing the domain,
        the grid, and a \texttt{RebarElementRange} that separates
        hex element indices $[0,\,n_{\mathrm{hex}})$ from truss indices
        $[n_{\mathrm{hex}},\,n_{\mathrm{total}})$.
\end{enumerate}

Because the truss elements share hex nodes, the total node count is
unchanged---only the element count grows.


% =========================================================================
\section{Reaction-Based Effective Modulus}
\label{sec:reinf-reaction-emod}

With \texttt{MultiElementPolicy}, the type-erased
\texttt{FEM\_Element} does not expose \texttt{material\_points()}, so
volume-averaged stress/strain extraction is not directly available.
Instead, \texttt{solve\_reinforced()} computes effective moduli from
boundary reactions.

After the incremental nonlinear solve converges at $p = 1$, internal
forces~$\mathbf{f}_\text{int}$ are reassembled from the converged state
vector.  The $z$-component of the resultant reaction at the max-$z$
face gives the axial force:
\begin{equation}
  F_z = \sum_{i \in \mathcal{N}_{\max z}} f_{\text{int},z}^{(i)},
  \qquad
  \sigma_{zz}^{\text{eff}} = \frac{F_z}{A},
  \qquad
  E_{\text{eff}} = \frac{\sigma_{zz}^{\text{eff}}}
                        {\varepsilon_{zz}^{\text{avg}}},
  \label{eq:reinf-eeff}
\end{equation}
where $A = w\!\times\!h$ is the cross-section area and
$\varepsilon_{zz}^{\text{avg}} = \bar{u}_z^{\max z} / L$ is the
average imposed strain.


% =========================================================================
\section{Bug Fix: State Vector Accumulation in Incremental Solve}
\label{sec:reinf-bugfix}

During the integration testing, the computed $E_{\text{eff}}$ was found
to be approximately $20{\times}$ higher than the expected value of
${\sim}28{,}500~\mathrm{MPa}$ for fc~$=30$ concrete.  Detailed diagnostics
revealed that the model state vector at constrained (Dirichlet) nodes
accumulated imposed displacements across all load increments, rather than
reflecting only the final converged state.

\subsection{Root Cause}

The \texttt{commit\_state()} method in \texttt{NonlinearAnalysis} updated the
model's state vector as:

\begin{lstlisting}[language=C++, basicstyle=\ttfamily\small]
DMGlobalToLocal(dm, U, INSERT_VALUES, model_->state_vector());
VecAXPY(model_->state_vector(), 1.0, model_->imposed_solution());
\end{lstlisting}

PETSc's \texttt{DMGlobalToLocal} with \texttt{INSERT\_VALUES} populates
free-DOF entries from the global vector~$U$ but \emph{does not touch}
constrained entries---they retain whatever value was already stored.
At each of the $N$ incremental steps, \texttt{VecAXPY} adds
$p_k \cdot u_{\text{imposed}}$ to the untouched constrained entries,
producing a cumulative sum:
\begin{equation}
  u_{\text{state}}^{(N)} \;=\;
    \sum_{k=1}^{N} p_k \cdot u_{\text{imposed}}
  \;=\; u_{\text{imposed}} \sum_{k=1}^{N} \frac{k}{N}
  \;=\; u_{\text{imposed}} \cdot \frac{N+1}{2}.
  \label{eq:accumulation}
\end{equation}

For $N = 20$ steps this gives a factor of $10.5{\times}$, exactly matching
the observed displacement amplification ($u_{\text{state}} = 0.00105$
vs.\ $u_{\text{imposed}} = 0.0001$).

\subsection{Fix}

The fix is a single-line addition of \texttt{VecSet(model\_->state\_vector(), 0.0)}
before the \texttt{DMGlobalToLocal} call:

\begin{lstlisting}[language=C++, basicstyle=\ttfamily\small]
VecSet(model_->state_vector(), 0.0);
DMGlobalToLocal(dm, U, INSERT_VALUES, model_->state_vector());
VecAXPY(model_->state_vector(), 1.0, model_->imposed_solution());
\end{lstlisting}

This ensures constrained entries start from zero at each commit,
so the final state is exactly $0 + u_{\text{imposed}}$ as intended.

The same defensive zero-fill was applied to
\texttt{DynamicAnalysis} and \texttt{LinearAnalysis} for consistency,
even though those code paths were not affected in practice (they
call the update only once, not in a loop).

\subsection{Scope of Impact}

The bug affected \emph{any} incremental displacement-controlled nonlinear
solve that:
\begin{itemize}
  \item Used more than one load step ($N > 1$), \textbf{and}
  \item Read the model's \texttt{state\_vector()} at constrained nodes
        after the solve (e.g., for reaction-force computation or
        post-processing).
\end{itemize}

The SNES convergence was \emph{not} affected---the residual evaluation
in \texttt{FormResidual} correctly zeroed its local vector before
\texttt{DMGlobalToLocal}.  Only the post-solve state vector used for
post-processing was corrupted.  This explains why the solver ``converged''
but the extracted reactions were wrong.

Note that the original \texttt{SubModelSolver::solve()} method extracted
effective moduli from volume-averaged Gauss-point stresses (not from
the state vector at constrained nodes), and therefore was not affected
by this bug.


% =========================================================================
\section{Validation}
\label{sec:reinf-validation}

The integration test \texttt{test\_reinforced\_submodel.cpp} exercises the
full pipeline with a $0.30 \times 0.30 \times 1.00$~m prismatic volume
($n_x\!=\!2$, $n_y\!=\!2$, $n_z\!=\!4$), $f_c'\!=\!30$~MPa concrete,
and an imposed axial strain of $\varepsilon = 10^{-4}$.

\paragraph{Unreinforced case.}
With empty rebar, \texttt{solve\_reinforced()} returns
$E_{\text{eff}} = 28{,}420$~MPa, matching the analytical value
$E_e \approx 28{,}513$~MPa from the Ko--Bathe elastic parameters
($K_e = 11{,}096$, $G_e = 13{,}304$,
$E_e = 9K_eG_e/(3K_e + G_e)$) to within $0.3\%$.

\paragraph{Reinforced case.}
A single $16$mm-equivalent rebar ($A_s = 0.0002$~m$^2$,
$\rho = A_s/A = 0.22\%$) with $E_s = 200{,}000$~MPa steel adds
$\Delta E \approx \rho \cdot E_s \approx 444$~MPa.  The measured
$E_{\text{eff}} = 28{,}865$~MPa gives $\Delta E = 445$~MPa, in excellent
agreement with the mixture rule.

All 14 assertions pass, verifying:
\begin{itemize}
  \item Domain construction with mixed hex/truss topology (7~checks),
  \item Convergence of both unreinforced and reinforced solves,
  \item Positive $E_{\text{eff}}$ in the physical range $[10{,}000,\; 50{,}000]$~MPa,
  \item Reinforced $E_{\text{eff}} > $ unreinforced (rebar adds stiffness),
  \item Reinforced $E_{\text{eff}} < 2\times$ unreinforced (small $\rho$).
\end{itemize}
